%============================================================
% CH2 · LESSON 4 — CYCLE DETECTION AND RECOVERY
%============================================================
\subsection{Cycle Detection and Recovering the Cycle}
\label{ssec:ch2-cycles}

``Does this graph have a cycle, and if so, show me one'' is the most common question DFS answers in contests: dependency graphs must be acyclic, a round trip must exist, a process must not loop forever. Lesson~\ref{ssec:ch2-dfs-tree} already did most of the work: back edges are the cycles.

\subsubsection*{The cycle theorem}

\begin{theorembox}[Back edges and cycles]
\begin{enumerate}[leftmargin=*]
    \item \textbf{Directed:} a directed graph has a cycle if and only if a DFS of it (from every vertex, as in the outer loop) finds a back edge --- an edge $u \to v$ to a \emph{gray} vertex $v$.
    \item \textbf{Undirected:} an undirected graph has a cycle if and only if a DFS finds an edge, other than the one it came in by, to an already visited vertex.
\end{enumerate}
In both cases the cycle is the back edge plus the tree path from $v$ down to $u$.
\end{theorembox}

\textbf{Proof (directed).} ($\Leftarrow$) A back edge $u \to v$ goes to a gray vertex, i.e.\ an ancestor of $u$; the tree path $v \leadsto u$ plus the edge $u \to v$ is a directed cycle.
($\Rightarrow$) Let $C$ be a directed cycle and $v$ the vertex of $C$ that DFS enters first. At time $\tin{v}$ all other vertices of $C$ are white, so the rest of the cycle is a white path from $v$ to its predecessor $u$ on $C$. By the white-path theorem $u$ becomes a descendant of $v$, and when DFS examines the edge $u \to v$, $v$ is still gray (it is left only after its descendants): a back edge. \hfill$\blacksquare$

The undirected case is the same argument with the extra care of Lesson~\ref{ssec:ch2-dfs-tree}: the edge to the parent is skipped, so the non-tree edges are exactly the back edges.

\begin{diagrambox}
\begin{tikzpicture}[x=0.95cm, y=0.95cm]
    \begin{scope}
        \tikzset{g2-v1/.style=gvsel, g2-v2/.style=gvsel, g2-v3/.style=gvcur,
                 g2-e12/.style={gtree, -{Stealth[length=2.4mm]}}, g2-e23/.style={gtree, -{Stealth[length=2.4mm]}},
                 g2-e31/.style=gback}
        \pic {g2-D};
        \node[font=\scriptsize, text width=6cm, align=left, anchor=north west] at (-0.4,-2.0)
            {The first back edge DFS meets on $D$ is $3 \to 1$ (at time 2, stack \texttt{1 2 3}). Walk the parents from 3 up to 1: $3, 2, 1$; reverse it and close the loop: \textbf{cycle $1 \to 2 \to 3 \to 1$}.};
    \end{scope}
    \begin{scope}[xshift=7.6cm]
        \tikzset{g2-u1/.style=gvsel, g2-u2/.style=gvsel, g2-u5/.style=gvsel, g2-u4/.style=gvcur,
                 g2-f12/.style=gtree, g2-f25/.style=gtree, g2-f45/.style=gtree, g2-f14/.style={gback, -}}
        \pic {g2-G3};
        \node[font=\scriptsize, text width=6cm, align=left, anchor=north west] at (-0.4,-2.0)
            {On $G_3$ the first non-parent edge to a visited vertex is $4$--$1$ (stack \texttt{1 2 5 4}). Parents from 4 up to 1: $4, 5, 2, 1$: \textbf{cycle $1, 2, 5, 4, 1$}.};
    \end{scope}
\end{tikzpicture}
\figcap{Recovering a cycle. Orange $=$ the vertex where the back edge (red, dashed) is found; purple $=$ the rest of the cycle, read off the parent pointers.}
\end{diagrambox}

\subsubsection*{Recovering the cycle}

When DFS sits at $u$ and finds the back edge $u \to v$, the cycle is on the call stack: it is the part of the stack from $v$ to $u$. Two ways to read it:
\begin{itemize}[leftmargin=*]
    \item \textbf{parent pointers:} $x = u$; while $x \ne v$: output $x$, $x = \texttt{parent}[x]$; output $v$. This lists the cycle backwards ($u, \dots, v$); reverse it for a directed cycle.
    \item \textbf{explicit path:} keep the stack as a vector \texttt{path}; the cycle is \texttt{path} from the position of $v$ to the end. (Problem~\ref{ssec:ch2-round-trip-ii} uses parents, Exercise E1.1 uses the path.)
\end{itemize}
Both cost \bigO{\text{cycle length}} once, after an \VE{} search.

\begin{lstlisting}[style=cpp]
#include <bits/stdc++.h>
using namespace std;

// Find ONE directed cycle with colours 0/1/2 and print it, or report that the graph is acyclic.
vector<vector<int>> adj;
vector<int> color, par;
int cycStart = -1, cycEnd = -1;                     // the back edge cycEnd -> cycStart

bool dfs(int u) {
    color[u] = 1;
    for (int v : adj[u]) {
        if (color[v] == 0) {
            par[v] = u;
            if (dfs(v)) return true;                // stop as soon as a cycle is known
        } else if (color[v] == 1) {                 // gray: a back edge u -> v
            cycStart = v; cycEnd = u;
            return true;
        }                                           // black: forward or cross edge, harmless
    }
    color[u] = 2;
    return false;
}

int main() {
    int n, m;
    cin >> n >> m;
    adj.assign(n + 1, vector<int>());
    for (int i = 0; i < m; i++) {
        int a, b;
        cin >> a >> b;
        adj[a].push_back(b);
    }
    color.assign(n + 1, 0); par.assign(n + 1, 0);
    for (int s = 1; s <= n && cycStart == -1; s++)
        if (color[s] == 0) dfs(s);
    if (cycStart == -1) { cout << "acyclic\n"; return 0; }
    vector<int> cyc;
    for (int x = cycEnd; x != cycStart; x = par[x]) cyc.push_back(x);
    cyc.push_back(cycStart);
    reverse(cyc.begin(), cyc.end());                // cycStart -> ... -> cycEnd
    cyc.push_back(cycStart);                        // close the loop
    for (int x : cyc) cout << x << " ";
    cout << "\n";
}
\end{lstlisting}

\begin{lstlisting}[style=trace]
input : D                                   output: 1 2 3 1
input : 3 3 / 1 2 / 1 3 / 2 3   (a DAG)     output: acyclic
input : 2 1 / 2 2               (self-loop) output: 2 2
\end{lstlisting}

\subsubsection*{Three traps}

\begin{diagrambox}
\begin{tikzpicture}[x=1cm, y=1cm]
    % trap 1: visited is not enough
    \begin{scope}
        \node[font=\small\bfseries] at (1,1.55) {(1) a black target};
        \node[gvdone] (a1) at (0,0.6) {1};
        \node[gvdone] (a2) at (2,0.6) {2};
        \node[gvcur] (a3) at (1,-0.6) {3};
        \draw[gtree, -{Stealth[length=2.4mm]}] (a1) -- (a2);
        \draw[gtree, -{Stealth[length=2.4mm]}] (a2) -- (a3);
        \draw[gfwd] (a1) -- (a3);
        \node[font=\scriptsize, text width=3.6cm, align=center] at (1,-1.6) {DAG $1{\to}2{\to}3$, $1{\to}3$: when 1 examines $1 \to 3$, vertex 3 is \textbf{black}. ``visited $\Rightarrow$ cycle'' would be wrong.};
    \end{scope}
    % trap 2: parallel edges
    \begin{scope}[xshift=5.2cm]
        \node[font=\small\bfseries] at (1,1.55) {(2) parallel edges};
        \node[gvseen] (b1) at (0,0) {1};
        \node[gvcur] (b2) at (2,0) {2};
        \draw[gtree] (b1) to[bend left=25] node[gw, above] {e0} (b2);
        \draw[gback, -] (b2) to[bend left=25] node[gw, below] {e1} (b1);
        \node[font=\scriptsize, text width=3.6cm, align=center] at (1,-1.6) {Undirected, two roads between 1 and 2: a cycle $1, 2, 1$ of length 2. ``Skip the parent \emph{vertex}'' misses it; skip the parent \emph{edge} (e0).};
    \end{scope}
    % trap 3: self loops
    \begin{scope}[xshift=10.4cm]
        \node[font=\small\bfseries] at (1,1.55) {(3) self-loops};
        \node[gvcur] (c1) at (1,0) {2};
        \draw[gback] (c1) to[out=60, in=120, looseness=7] (c1);
        \node[font=\scriptsize, text width=3.6cm, align=center] at (1,-1.6) {$2 \to 2$ is a back edge to itself: a cycle of length 1. Fine if the problem allows it --- many statements (Round Trip) forbid short cycles by guaranteeing simple graphs.};
    \end{scope}
\end{tikzpicture}
\figcap{The three classic situations a cycle detector must handle correctly.}
\end{diagrambox}

\begin{warningbox}[Directed graphs: gray, not visited]
The test must be ``$v$ is gray'' (on the current path). A black vertex was finished by an earlier part of the search; reaching it again only means two paths lead there (trap 1). With the test changed to ``$v$ is not white'', the program above declares a cycle at the edge $1 \to 3$ of the 3-vertex DAG in the trace --- and then the parent walk from 1 never meets 3, because 3 is not an ancestor of 1: our run looped until it died with \texttt{bad\_alloc}. A wrong cycle test often shows up as a crash in the \emph{recovery} code, not as a wrong answer.
\end{warningbox}

\begin{warningbox}[Undirected graphs: skip the edge, not the vertex]
In an undirected multigraph, two parallel edges $e_0, e_1$ between $p$ and $v$ form a cycle. When $v$ was entered through $e_0$, skip only $e_0$; then $e_1$ is a non-tree edge to the visited $p$ and the cycle $p, v, p$ is found. If the input guarantees a simple graph (no parallel edges), skipping the parent vertex is equivalent and slightly shorter.
\end{warningbox}

\subsubsection*{Undirected graphs: counting instead of searching}

For the \emph{yes/no} question in an undirected graph, a search is not even needed: a forest with $n$ vertices and $c$ components has exactly $n - c$ edges, so
\[
    \text{the graph has a cycle} \iff m > n - c.
\]
Counting the components $c$ still needs a traversal (or a DSU, Chapter~6), but this formula is handy for checks like ``is the graph a tree?'' ($m = n - 1$ and connected) or ``does every component have exactly one cycle?'' ($m = n$ per component, Exercise E1.2). To \emph{print} a cycle you need the DFS.

\subsubsection*{Iterative versions}

On $10^5$ vertices a long path can make the recursion deep (Lesson~\ref{ssec:ch2-call-stack}). Both detectors run unchanged on the iterator-based iterative DFS: colour 1 when a vertex is pushed, colour 2 when it is popped, and check the colour of every examined neighbour. The library file \texttt{code/lib/cycle.hpp} (Appendix~A) contains both iterative versions, the undirected one with edge ids; Problems~\ref{ssec:ch2-round-trip} and~\ref{ssec:ch2-round-trip-ii} show them on real inputs.

\subsubsection*{Check your understanding}

\begin{enumerate}[leftmargin=*]
    \item In a directed graph DFS finds the back edge $7 \to 6$ of $D$ (if it had not stopped at $3 \to 1$). Which cycle does the parent walk print?
    \item True or false: if a directed graph has a cycle, DFS from vertex 1 alone finds it.
    \item An undirected graph has $n = 10$ vertices, $m = 9$ edges and 2 components. Does it have a cycle?
    \item Why is it safe to stop the whole search at the first back edge, even though the colours of the unfinished vertices are then wrong?
    \item A directed graph has $10^5$ vertices and $2\cdot10^5$ edges and you must print a cycle. Which of the three traps can appear, and what is the time complexity?
\end{enumerate}

\begin{solutionbox}[Answers]
\begin{enumerate}[leftmargin=*]
    \item $\texttt{par}[7] = 6$, so the walk is $7, 6$; reversed and closed: $6 \to 7 \to 6$.
    \item False: the cycle may not be reachable from 1 (in $D$, a cycle through 8 would be invisible from 1). That is why the outer loop starts a DFS from every still-white vertex.
    \item A forest with 10 vertices and 2 components has $10 - 2 = 8$ edges; $9 > 8$, so \textbf{yes}.
    \item We need one cycle, and it is fully known at that moment (the back edge plus the parent pointers, which are final for all gray vertices). The colours are never used again.
    \item Trap 1 (black targets) and trap 3 (self-loops, if allowed). Trap 2 does not exist in directed graphs: $u \to v$ and $v \to u$ are two different edges and genuinely form a cycle of length 2. Time \bigO{n + m}.
\end{enumerate}
\end{solutionbox}
